\documentclass[10pt]{article}
\usepackage[T1]{fontenc}
\usepackage[utf8]{inputenc}
\usepackage[french]{babel}
\usepackage{amsmath,amssymb,mathrsfs,mathpazo}
\usepackage{geometry}
\geometry{a4paper,margin=2cm}
\usepackage{graphicx,array,tabularx,tabularray,colortbl,multirow,multicol,booktabs,enumitem,pifont,ulem,fancybox,diagbox,cancel}
\usepackage{pstricks,pstricks-add,pst-plot,pst-tree,pst-node,pst-text,pst-func,pst-math,pst-bspline,pst-3dplot,pst-eucl,pst-solides3d}
\let\pscalculate\undefined
\usepackage{pst-bezier}
\usepackage{tikz,pgfplots,tkz-tab}
\usetikzlibrary{arrows,arrows.meta,patterns,calc,positioning,trees,decorations.markings,shapes.geometric,shapes.misc,shapes.arrows,intersections,decorations.pathmorphing,decorations.pathreplacing,backgrounds,shadows,babel}
\pgfplotsset{compat=1.18}
\usepgfplotslibrary{fillbetween}
\usepackage[np]{numprint}
\usepackage{listings,makecell,pas-tableur,twoopt}
\let\barre\undefined
\usepackage[pstricks]{tablvar}
\definecolor{darkgreen}{RGB}{0,100,0}
\providecommand{\red}{\color{red}}
\providecommand{\blue}{\color{blue}}
\providecommand{\green}{\color{green}}
\providecommand{\black}{\color{black}}
\providecommand{\white}{\color{white}}
\providecommand{\gray}{\color{gray}}
\providecommand{\cyan}{\color{cyan}}
\providecommand{\magenta}{\color{magenta}}
\providecommand{\yellow}{\color{yellow}}
\providecommand{\violet}{\color{violet}}
\providecommand{\orange}{\color{orange}}
\providecommand{\euro}{\texteuro}
\providecommand{\EUR}{\texteuro}
\providecommand{\ieme}{\textsuperscript{e}}
\providecommand{\degres}{\ensuremath{^\circ}}
\providecommand{\degree}{\ensuremath{^\circ}}
\providecommand{\vv}[1]{\overrightarrow{#1}}
\providecommand{\Coord}[1]{#1}
\providecommand{\decosix}{}
\providecommand{\starredbullet}{\ensuremath{\bullet}}
\pagestyle{empty}
\setlength{\parindent}{0pt}
\usepackage{adjustbox}\DeclareUnicodeCharacter{2236}{:}\DeclareUnicodeCharacter{2190}{\ensuremath{\leftarrow}}
\providecommand{\R}{\mathbb{R}}
\providecommand{\N}{\mathbb{N}}
\providecommand{\Z}{\mathbb{Z}}
\providecommand{\Q}{\mathbb{Q}}
\providecommand{\D}{\mathbb{D}}
\providecommand{\C}{\mathbb{C}}
\providecommand{\e}{{\mathrm{e}}}
\providecommand{\vect}[1]{\overrightarrow{#1}}
\providecommand{\vectt}[1]{\overrightarrow{\mathrm{#1}}}
\providecommand{\vvt}[1]{\overrightarrow{\mathrm{#1}}}
\providecommand{\np}[2][]{{#2~#1}}
\providecommand{\vv}[1]{{\overrightarrow{#1}}}
\renewcommand{\d}{{\mathrm{d}}}
\providecommand{\degres}{{^\circ}}
\providecommand{\degre}{{^\circ}}
\providecommand{\blue}{\color{blue}}
\providecommand{\red}{\color{red}}
\providecommand{\enskip}{\quad}
\providecommand{\llbracket}{\lbrack\!\lbrack}
\providecommand{\rrbracket}{\rbrack\!\rbrack}
\providecommand{\euro}{{\text{€}}}
\providecommand{\no}{{\text{n°}}}
\providecommand{\interval}[2]{{[#1;#2]}}
\providecommand{\intervFO}[2]{{[#1;#2[}}
\providecommand{\dotfill}{\ldots}

\providecommand{\ds}{\displaystyle}
\providecommand{\textding}[1]{#1}
\providecommand{\Oij}{\(\mathrm{O};\vec{\imath},\vec{\jmath}\)}
\providecommand{\Oijk}{\(\mathrm{O};\vec{\imath},\vec{\jmath},\vec{k}\)}
\providecommand{\Ouv}{\(\mathrm{O};\vec{u},\vec{v}\)}
\providecommand{\intervFF}[2]{\ensuremath{[#1;#2]}}
\renewcommand{\euro}{\ifmmode\text{\texteuro}\else\texteuro\fi}
\DeclareUnicodeCharacter{00B0}{\ifmmode^{\circ}\else\textdegree\fi}
\DeclareUnicodeCharacter{20AC}{\ifmmode\text{\texteuro}\else\texteuro\fi}
% Énoncés : Éducation nationale. Transcription et corrigés : APMEP, Denis Vergès et contributeurs. Découpage et classement : CoopMaths, Sébastien Lozano, Éric Elter et contributeurs. Adaptation de la présentation : Maths Explorer.
% CC BY-SA : https://coopmaths.fr/www/about/
% https://coopmaths.fr/alea/static/eam/2026/tex/eam_specifique_2026_04_polynesie_3.tex
\begin{document}
\noindent\textit{APMEP / CoopMaths. CC BY-SA. Présentation adaptée par Maths Explorer.}\par\medskip
Le nombre hebdomadaire de consultations lors d'une épidémie de grippe saisonnière, pendant 10 semaines consécutives, est modélisé par la fonction $f$, définie sur l'intervalle $[0 ; 10]$ par :

$f(t)=-25t^{2}+300t+400$ où $t$ représente le nombre de semaines écoulées depuis le début de l'épidémie.

La courbe représentative $\mathcal{C}_{f}$ de la fonction $f$ est donnée ci-dessous dans un repère du plan.

\begin{figure}
	\begin{center}
		\begin{tikzpicture}[x=1.2cm, y=0.0075cm]
			
			\draw[gray!30, line width=0.3pt] (0,0) grid[xstep=0.2, ystep=20] (10.2,1420);
			
			\draw[gray!60, line width=0.6pt] (0,0) grid[xstep=1, ystep=100] (10.2,1420);

			
			\draw[line width=1pt, ->, >=stealth] (0,0) -- (10.6,0) node[below left, font=\small\bfseries] at (10,-50){nombre de semaines};
			\draw[line width=1pt, ->, >=stealth] (0,0) -- (0,1450) node[below right, font=\small\bfseries] at (0.1,1410){Nombre de consultations};

			
			\foreach \x in {0,1,...,10}
			\draw (\x,10) -- (\x,-10) node[below, font=\small]{$\np{\x}$};

			
			\foreach \y in {0,100,...,1400}
			\draw (0.08,\y) -- (-0.08,\y) node[left, font=\small]{$\np{\y}$};

			
			\draw[red!70!black, line width=1.2pt] plot[smooth, domain=0:10, samples=100] (\x, {-25*\x*\x + 300*\x + 400});
			\node[above right, font=\large] at (8.8, 920) {$\mathcal{C}_f$};
		\end{tikzpicture}
	\end{center}
\end{figure}

\subsection*{Partie A : lecture graphique}

En utilisant le graphique ci-dessus :
\begin{enumerate}
	\item Déterminer le nombre de consultations 2 semaines après le début de l'épidémie.
	\item Déterminer l'intervalle de temps pour lequel le nombre de consultations dépasse \np{1200}.
	\item Au bout de combien de semaines après le début de l'épidémie le nombre de consultations est-il maximal ? Quelle est cette valeur maximale ?
\end{enumerate}

\subsection*{Partie B : Étude de la fonction $f$}

\begin{enumerate}
	\item Calculer $f(0)$. Interpréter le résultat dans le contexte de l'exercice.
	\item Justifier que, pour tout réel $t$ de l'intervalle $[0~;~10]$, $f'(t) = -50t + 300$.
	\item Résoudre dans l'intervalle $[0~;~10]$ l'équation $f'(t)=0$.
	\item Étudier le signe de $f'(t)$, puis dresser le tableau de variation de la fonction $f$ sur l'intervalle $[0~;~10]$. On donne $f(6)=\np{1300}$.
\end{enumerate}

\end{document}
